\section{Residues, stability, and local coefficients}\label{sec:residues}

\subsection{Finite schemes and the normal-form residue}

\begin{lemma}\label{lem:normal-form}
For every $n\ge1$ and every parameter choice, $Z_n$ is a finite complete
intersection of length $d^n$, with no common zero at infinity in the
homogenized cyclic equations. For any polynomial $g(\xx)$,
\begin{equation}\label{eq:normal-form}
 \Res_{F_n}(g)=
 [x_0^{d-1}\cdots x_{n-1}^{d-1}]\NF_{F_n}(g).
\end{equation}
In particular, the residue of $g_n$ is a polynomial in the parameters.
\end{lemma}

\begin{proof}
Choose any monomial order that first compares total degree. Since
$d\ge2$, the leading monomial of $F_{n,i}$ is $x_i^d$, including the
two short-period cases \eqref{eq:period-one}--\eqref{eq:period-two}.
These leading monomials are pairwise relatively prime. The relatively
prime leading-monomial criterion therefore shows that the $F_{n,i}$
form a Gr\"obner basis. The standard monomials are exactly
\[
 x_0^{j_0}\cdots x_{n-1}^{j_{n-1}},\qquad 0\le j_i<d.
\]
They give a basis of the quotient algebra of size $d^n$. An ideal of
finite colength generated by $n$ polynomials in $n$ variables is a
zero-dimensional complete intersection. Its scheme length is the
dimension just computed.

The top homogeneous equations are $x_i^d=0$ for all $i$. Their only
common affine zero is the origin, so they have no common projective
zero. This verifies the no-infinity hypothesis for the global residue.
The pure-power initial-form residue formula
\citep[Lemma~4.2 of the author preprint]{cattani1996residues} gives
\eqref{eq:normal-form}. The formula includes multiplicities through
local residues and requires no reducedness assumption.

All reductions use monic polynomials, so the division procedure never
divides by a parameter or a discriminant. Applying
\eqref{eq:normal-form} to $g_n$ therefore gives a polynomial in the
$p_j$, $a$, and $q_j$ even at parameters where the fixed scheme is
nonreduced.
\end{proof}

\subsection{The sign in the return denominator}

At a cyclic point, put
\begin{equation}\label{eq:monodromy}
 v_i=p'(x_i),\qquad
 A_i=\begin{pmatrix}v_i&-a\\1&0\end{pmatrix},\qquad
 \mathcal M_n=A_{n-1}\cdots A_0.
\end{equation}
The chronological order in \eqref{eq:monodromy} gives
$\mathcal M_n=DH^n(x_0,x_{n-1})$.

\begin{lemma}[Cyclic Hill sign]\label{lem:hill}
Let $J_n=DF_n$ be the cyclic Jacobian, with all coincident incidences
added. Then
\begin{equation}\label{eq:hill}
 \det J_n=\Tr(\mathcal M_n)-1-a^n
          =-\det(I-\mathcal M_n)
\end{equation}
for every $n\ge1$. Hence the residue trace
\eqref{eq:residue-trace} agrees with the pointwise sum
\eqref{eq:pointwise-trace} on the nondegenerate locus.
\end{lemma}

\begin{proof}
For $n\ge3$, the cyclic matrix $J_n$ has diagonal $v_i$, right-neighbour
entries $-1$, and left-neighbour entries $-a$. In the scalar cyclic
formula of \citet[Eq.~(1) of the author preprint v3]{molinari2008determinants},
substitute these entries. The two corner products contribute
\[
 (-1)^{n+1}\bigl((-1)^n+(-a)^n\bigr)=-1-a^n,
\]
and the transfer-product term is $\Tr(\mathcal M_n)$. Its nonzero
off-diagonal assumptions hold since $a\ne0$. This gives the first
equality in \eqref{eq:hill}. Also $\det A_i=a$, so
$\det\mathcal M_n=a^n$. The identity
$\det(I-M)=1-\Tr M+\det M$ for a two-by-two matrix proves the second
equality.

For $n=1$, direct differentiation gives
$\det J_1=v_0-1-a$, while $\Tr\mathcal M_1=v_0$.
For $n=2$, the neighbours coincide and
\[
 J_2=\begin{pmatrix}v_0&-(1+a)\\-(1+a)&v_1\end{pmatrix},
 \qquad \Tr(A_1A_0)=v_0v_1-2a.
\]
Thus
\[
 \det J_2=v_0v_1-(1+a)^2
          =\Tr(A_1A_0)-1-a^2.
\]
The sign is therefore independent of $n$.

At a simple common zero, the local residue of $g$ is
$g/\det J_n$. By \eqref{eq:hill}, its negative equals
$g/\det(I-DH^n)$. Summing over the fixed points and using
\eqref{eq:scheme-correspondence} proves the final assertion.
\end{proof}

\subsection{A local exponent expansion}

The next formula separates the single-coordinate residue coefficients
from the nearest-neighbour combinatorics. Choose a radius $R$ large
enough that
\begin{equation}\label{eq:radius}
 |p(X)|>(1+|a|)R\qquad\text{when }|X|=R.
\end{equation}
Such a choice exists since $d\ge2$. It depends on $p,a$, but not on $n$.

\begin{lemma}\label{lem:local-expansion}
For each $n\ge1$, the global residue $\rho_n$ is the convergent sum
\begin{equation}\label{eq:configuration-sum}
 \rho_n=
 \sum_{\substack{L_0,R_0,\ldots,L_{n-1},R_{n-1}\ge0}}
 \prod_{i=0}^{n-1}
 \binom{L_i+R_i}{L_i}a^{L_i}
 c_{R_{i-1}+L_{i+1},\,L_i+R_i}(p,q).
\end{equation}
If a product in this sum is nonzero, then with $e_i=L_i+R_i$ it satisfies
\begin{equation}\label{eq:flow-inequality}
 d e_i\le m-d+1+R_{i-1}+L_{i+1}
 \qquad(0\le i<n).
\end{equation}
\end{lemma}

\begin{proof}
The global coefficient formula for a system with these pure-power
initial forms \citep[Theorem~2.3 of the author preprint]{cattani1996residues}
identifies its residue with the coefficient of
$x_0^{-1}\cdots x_{n-1}^{-1}$ in the expansion at infinity. Equivalently,
on a sufficiently large product torus it is
\begin{equation}\label{eq:torus-residue}
 \rho_n=\frac{1}{(2\pi i)^n}
 \int_{|x_0|=\cdots=|x_{n-1}|=R}
 \frac{g_n(\xx)\,dx_0\cdots dx_{n-1}}
      {\prod_{i=0}^{n-1}F_{n,i}(\xx)},
\end{equation}
where each circle has positive orientation and the factor order is as
in Definition~\ref{def:residue-trace}. In this pure-power setting the
large-torus formula is the coefficient formula just cited. The lower
degree terms may be expanded as convergent geometric series at large
$R$, and regrouping the terms involving only $x_i$ gives
\begin{equation}\label{eq:geometric-expansion}
 \frac{1}{F_{n,i}}=
 \sum_{L_i,R_i\ge0}\binom{L_i+R_i}{L_i}
 \frac{a^{L_i}x_{i-1}^{L_i}x_{i+1}^{R_i}}
      {p(x_i)^{L_i+R_i+1}}.
\end{equation}
Condition~\eqref{eq:radius} makes this expansion uniformly absolutely
convergent on the integration torus. For each fixed $n$, its finite
product can therefore be integrated term by term. Uniformity of the
whole product norm in $n$ is not needed.

The variable $x_i$ receives $R_{i-1}$ powers from its left neighbour
and $L_{i+1}$ powers from its right neighbour. Its one-variable integral
is consequently $c_{R_{i-1}+L_{i+1},L_i+R_i}$. Multiplying the local
factors proves \eqref{eq:configuration-sum}.

For $n=1$, both neighbouring factors use the same variable. Summing
over $L_0+R_0=e$ produces the coefficient
$\sum_{L_0=0}^e\binom e{L_0}a^{L_0}=(1+a)^e$, as required by
\eqref{eq:period-one}. For $n=2$, the left and right neighbours again
coincide, but their exponents add in the same variable. Formula
\eqref{eq:configuration-sum} therefore retains the correct two
incidences without treating them as different coordinates.

Finally, if $c_{k,e}\ne0$, the highest possible power at infinity in
\eqref{eq:local-coefficient} cannot be below $X^{-1}$. Hence
\begin{equation}\label{eq:local-degree-condition}
 m+k-d(e+1)\ge-1.
\end{equation}
Apply this necessary condition to each local factor of a nonzero
product to obtain \eqref{eq:flow-inequality}. No sufficiency claim is
made for this degree condition.
\end{proof}
